% ================= CHAPTER 3 — THEORETICAL BACKGROUND AND TECHNOLOGIES =================
% Length: <= 10 pages. Summarise existing knowledge; for EACH technology, say which
% Chapter-2 requirement it solves, list alternatives, and justify the choice. Cite sources.
This chapter introduces the theoretical foundations and technologies used in the
system. For each technology, the requirement it addresses, the available
alternatives, and the rationale for the selection are discussed.

\section{Inventory Management Theory}
% GUIDE: stock control, reservation, weighted-average costing, and the FEFO
% (First-Expired-First-Out) strategy vs FIFO/LIFO for perishable goods. Cite.
Inventory management controls the quantity, location, cost, and movement history
of goods so that business operations can meet demand without creating excessive
stock. In a distribution company, the inventory record is not only an accounting
figure but also an operational constraint: a sales order should be accepted only
when the requested product can be supplied from a warehouse, and a warehouse
transaction should leave an audit trail from the business document that caused
it. Therefore, the system distinguishes three quantities for each product and
warehouse: on-hand quantity, reserved quantity, and available quantity. On-hand
quantity represents the physical stock in the warehouse. Reserved quantity
represents stock committed to approved sales orders but not yet shipped.
Available quantity is calculated as on-hand minus reserved stock and is used to
prevent overselling.

For goods with expiry dates, stock control must also identify the lot from which
each outbound quantity is taken. General warehouse-management literature
emphasizes that lot traceability and stock rotation are central to reducing loss
and improving warehouse discipline \cite{richards2021warehouse}. FIFO
(First-In-First-Out) issues the oldest received stock first, while LIFO
(Last-In-First-Out) issues the newest received stock first and is mainly an
inventory valuation convention. In contrast, FEFO (First-Expired-First-Out)
issues the stock with the earliest expiry date first. The management of
perishable stock, whose value falls as the expiry date approaches, has long been
studied in inventory theory~\cite{nahmias1982perishable}. FEFO is more
appropriate for food, medicine, cosmetics, and other perishable goods because the
receiving date and the expiry date are not always aligned. Odoo's inventory documentation,
for example, separates expiration-date tracking from the FEFO removal strategy,
which reflects the same principle \cite{odooExpiration,odooFefo}. This thesis
therefore models both aggregate inventory and per-lot inventory, then performs
lot allocation during goods issue.

The system uses moving weighted-average cost for the aggregate inventory record.
When goods are received, the previous stock value and the new receipt value are
combined to update the average cost. This approach is suitable for the project
scope because it avoids the complexity of full accounting valuation by batch
while still keeping cost information for stock transactions. The FEFO mechanism
is used for physical issue sequence, whereas the weighted average is used for
simple inventory-cost reporting.

\section{Backend Technologies}
% GUIDE: Java 17, Spring Boot (IoC, layered architecture), Spring Data JPA/Hibernate
% (ORM, optimistic locking via @Version), Spring Security + JWT (stateless RBAC),
% Flyway (schema migration). For each: requirement solved + alternatives + why.
The backend is a server-side application that exposes the business functions
through REST APIs. This section describes each backend technology in turn: the
requirement it addresses and the reason it was selected over the alternatives.

\subsection{Java and Spring Boot}
The backend is implemented in Java 17 with Spring Boot 3.1.4
\cite{springboot314}. Spring Boot was selected because it provides an integrated
platform for building standalone REST applications and has explicit support for
Java 17 in version 3.1.x. The requirement it addresses is a maintainable
server-side application that exposes business functions through stable APIs.
Node.js/Express and Django were considered as alternatives, but Spring Boot fits
the project's domain better because the business rules are transaction-heavy and
map naturally to service classes, repositories, and declarative transaction
management. On top of Spring Boot, the backend follows a layered architecture, an
established enterprise pattern that separates presentation, domain logic, and
data access~\cite{fowler2002poeaa}: controllers receive HTTP requests and
validate input DTOs; services implement business rules such as purchase approval,
goods-receipt confirmation, inventory reservation, FEFO allocation, and
delivery-status recovery; and repositories encapsulate database access. This
keeps each module's responsibility clear and testable at the service layer
without exposing persistence details to the frontend.

\subsection{Spring Data JPA and Hibernate}
Spring Data JPA, with Hibernate as the implementation, provides
object-relational mapping \cite{bauer2015hibernate}. It reduces manual SQL for
common CRUD operations and lets domain relationships such as
SalesOrder--SalesOrderItem, GoodsReceipt--GoodsReceiptItem, and
Inventory--Warehouse be expressed as Java entities. Plain JDBC or a SQL mapper
such as MyBatis were the alternatives: plain SQL offers fine-grained control but
produces repetitive code and makes entity relationships harder to maintain. JPA
was selected because most project operations are entity-centric, while custom
repository queries remain available for cases such as FEFO lot selection and
pessimistic locking during stock updates.

\subsection{Database Schema Management}
The database schema and seed data are defined as versioned SQL migration files
named after the Flyway naming convention (\texttt{V1}, \texttt{V2}, \dots) under
\texttt{db/migration} \cite{flywayMigrations}. The application itself does not
embed a migration runtime; instead, the scripts are applied to PostgreSQL in
version order with the \texttt{psql} client when an environment is prepared.
Hibernate runs with \texttt{ddl-auto=validate}, so it checks the entity mappings
against the existing schema rather than generating it. This file-based,
version-ordered scheme was selected because database changes must be repeatable
and reproducible across development and demonstration environments rather than
applied as ad-hoc edits to each database.

\section{Frontend Technologies}
% GUIDE: React 18 (component model, SPA), Vite, Ant Design, Axios, Recharts.
% Alternatives: Angular/Vue; explain the choice.
The frontend is implemented as a single-page application. This section describes
each frontend technology in turn: the requirement it addresses and the reason it
was selected over the available alternatives.

\subsection{React}
React 18 is the core user-interface library \cite{react18}. Its component model
supports decomposing the interface into reusable pages, forms, tables, and status
widgets, which the project uses for modules such as purchase orders, sales orders,
goods receipts, goods issues, inventory lots, delivery plans, and dashboard
reports. Angular and Vue were considered as alternatives: Angular provides a
complete framework but is heavier than this project requires, while Vue is
lightweight but has a smaller enterprise-component ecosystem. React was chosen
because its component model, combined with the mature React-based component
library described below, allows consistent business screens to be built quickly.

\subsection{Vite}
Vite 5 is the frontend build tool and development server \cite{vite5}. Compared
with older setups such as Create React App or a hand-written Webpack
configuration, Vite provides a faster development server and a simpler production
build, which is valuable during the frequent interface iteration of this project.
It addresses the requirement for an efficient and reproducible frontend build.

\subsection{Ant Design}
Ant Design 5 is the user-interface component library \cite{antdesign}. It supplies
enterprise-oriented components such as tables, forms, modals, date pickers, steps,
tags, and layout primitives. These match a distribution-management application, in
which users repeatedly scan tables, edit forms, approve documents, and monitor
statuses. Adopting a single, consistent component library also enforces a uniform
look across all modules without designing each control from scratch.

\subsection{React Router}
React Router 6 provides client-side routing \cite{reactRouter}. It maps browser
URLs to the application's pages and nested layouts so that navigation between the
dashboard, document lists, and detail screens happens within the single-page
application without full page reloads. This supports the requirement for a
multi-screen interface with clear navigation.

\subsection{Axios}
Axios is the HTTP client used to communicate with the REST backend \cite{axios}.
It centralizes API calls, attaches the JWT access token to each request, and
serializes JSON request and response bodies. A shared Axios instance keeps
authentication and error handling consistent across all modules.

\subsection{Recharts}
Recharts is a React charting library used for the management dashboard
\cite{recharts}. It renders aggregate indicators, such as revenue and
top-selling-product reports, as bar and line charts. It was selected because it
integrates directly with the React component model used by the rest of the
frontend, avoiding a separate charting runtime.

\section{Database — PostgreSQL}
% GUIDE: relational model, ACID transactions, why PostgreSQL over MySQL/NoSQL here.
PostgreSQL is selected as the relational database management system. The system
data is highly relational: orders contain items, goods receipts refer to purchase
orders, goods issues refer to sales orders, invoices refer to goods issues,
delivery trips contain delivery orders, and inventory records are linked to
products and warehouses. A relational model is therefore appropriate because it
supports foreign keys, joins, constraints, and transactions. PostgreSQL also
provides mature SQL support and official long-term documentation
\cite{postgresql14}.

The most important database requirement is transactional consistency. For
example, confirming a goods receipt must update the receipt status, purchase
order received quantities, aggregate inventory, inventory transactions, and
inventory lots as one logical unit. Confirming a goods issue must deduct lots,
deduct aggregate inventory, release reservation, update delivered quantities,
create an invoice, and update sales order status. A document database could store
large nested documents, but it would make these cross-entity consistency rules
and reporting queries more difficult. PostgreSQL was therefore chosen over a
NoSQL database for integrity, queryability, and compatibility with JPA.

\section{Authentication and Authorization}
% GUIDE: JWT structure, stateless auth, BCrypt hashing, role-based access control.
Authentication is based on JSON Web Tokens (JWT). According to RFC 7519, a JWT is
a compact representation of claims transferred between parties \cite{rfc7519}.
In this system, after a successful login the server signs an access token that
contains the username, user identifier, full name, and roles. The frontend sends
the token in subsequent API requests, and a Spring Security filter validates the
signature and expiration before allowing access to protected endpoints. This
stateless design avoids storing server-side HTTP sessions and is suitable for a
React SPA communicating with a REST backend.

Passwords are protected with BCrypt through Spring Security's
\texttt{Password\allowbreak Encoder}. Authorization follows Role-Based Access Control
(RBAC). The NIST RBAC model formalizes the mapping among users, roles, and
permissions \cite{sandhu2000nist}. The project applies this idea through nine
business roles: administrator, purchase staff, purchase manager, sales staff,
sales manager, warehouse staff, delivery administrator, shipper, and accountant.
Endpoint-level rules restrict which roles can create, update, approve, confirm,
or delete documents. Service-level rules add domain restrictions, such as
preventing warehouse staff from seeing unapproved orders and allowing shippers to
work only with assigned delivery trips. This combination supports separation of
duties in the procure-to-stock and order-to-cash workflows.
